\documentclass[a4paper,12pt]{report}
\usepackage[russian,custom]{config}

\begin{document}
\chapter{Ссылки}\label{chap:main}
Формы для рисунка: \объект{fig:main}; \объекта{fig:main}; \объекту{fig:main}; \объектВинительный{fig:main}; \объектом{fig:main}; \объекте{fig:main}.
Формы множественного числа для рисунка: \объекты{fig:main}; \объектов{fig:main}; \объектам{fig:main}; \объектыВинительный{fig:main}; \объектами{fig:main}; \объектах{fig:main}.
Список рисунков: \Объекты{fig:main}{fig:second}{fig:third}.
Список рисунков с продолжением: \Объекты{fig:main}{fig:second} демонстрируют пробел.
Формы для таблицы: \Объект{tab:main}; \объектВинительный{tab:main}; \объекте{tab:main}.
Формы множественного числа для таблицы: \Объекты{tab:main}; \объектов{tab:main}; \объектыВинительный{tab:main}; \объектах{tab:main}.
Формы для кода: \Объект{code:main}; \объектом{code:main}.
Формы множественного числа для кода: \Объекты{code:main}; \объектов{code:main}; \объектами{code:main}.
Формы для уравнения и главы: \Объект{eq:main}; \Объект{chap:main}.

\insertImage[fig:main]{cat_1.jpg}{Основной рисунок}
\insertImage[fig:second]{cat_1.jpg}{Второй рисунок}
\insertImage[fig:third]{cat_1.jpg}{Третий рисунок}
\insertTable[tab:main]{Основная таблица}{%
  \begin{tabular}{ll}
    A & B \\
  \end{tabular}
}
\begin{code}[code:main]{text}{Основной листинг}
main code
\end{code}
\begin{equation}
  \label{eq:main}
  a = b
\end{equation}
\end{document}
